% =============================================================================
% APPENDIX HI: CORE-HIERARCHY: Decision Stratification
% =============================================================================
% Module: BCM2_14_HIE (Decision Hierarchy)
% Category: CORE
% Version: 1.0 (January 2026)
% Status: Final
% Language: English
%
% PURPOSE: This appendix answers the question "At what level of the decision
% hierarchy does this choice occur?" It formalizes the four decision levels:
% L0 (Automatic), L1 (Heuristic), L2 (Deliberative), L3 (Meta-Cognitive).
%
% 10C POSITION: This is CORE 9 of 10 - the "HIERARCHY" question.
%
% =============================================================================

\begin{tcolorbox}[colback=gray!5!white, colframe=gray!75!black,
    title=Appendix Header: HI CORE-HIERARCHY]

\begin{tabular}{@{}ll@{}}
\textbf{Code:} & HI \\
\textbf{Category:} & CORE (Fundamental Theory) \\
\textbf{Full Name:} & CORE-HIERARCHY: Decision Stratification \\
\textbf{Module:} & BCM2\_14\_HIE (Decision Hierarchy) \\
\textbf{Version:} & 1.0 (January 2026) \\
\textbf{Status:} & Final \\
\end{tabular}

\vspace{0.5em}
\textbf{Purpose:} Answers the ninth fundamental question:
\textit{``At what level of the decision hierarchy does this choice occur?''}
Formalizes the four decision levels that determine processing mode and intervention type.

\end{tcolorbox}

\begin{tcolorbox}[colback=purple!5!white, colframe=purple!75!black,
    title=Chapter Linkage]

\textbf{Primary Chapter:} Chapter 13: Decision Architecture

\textbf{The Fundamental Question:}
\begin{quote}
\textit{``At what cognitive level is this decision being made? Is it automatic,
heuristic, deliberative, or meta-cognitive?''}
\end{quote}

\textbf{Relationship to Other COREs:}
\begin{itemize}[nosep]
\item Distinct from: Welfare Hierarchy (WHO) which concerns utility bearers
\item Affects: Processing mode, intervention appropriateness
\item Modulated by: Context (CORE-WHEN), cognitive load
\end{itemize}

\end{tcolorbox}

\chapter{CORE-HIERARCHY: Decision Stratification}
\label{app:core-hierarchy}

\begin{abstract}
This appendix answers the ninth fundamental question:
\textit{``At what level of the decision hierarchy does this choice occur?''}
We formalize four decision levels: $L_0$ (Automatic/System 1), $L_1$ (Heuristic),
$L_2$ (Deliberative/System 2), and $L_3$ (Meta-Cognitive). The level determines
which behavioral biases apply and which intervention types are effective.
This is distinct from the welfare hierarchy (CORE-WHO), which concerns who has utility.
\end{abstract}

\begin{tcolorbox}[colback=blue!5!white, colframe=blue!75!black,
    title=Quick Reference: CORE-HIERARCHY]

\textbf{Core Question:} At what cognitive level is the decision being made?

\textbf{The Four Decision Levels:}
\begin{itemize}[nosep]
\item $L_0$: \textbf{Automatic} --- Fast, unconscious, effortless (System 1)
\item $L_1$: \textbf{Heuristic} --- Rule-based, satisficing
\item $L_2$: \textbf{Deliberative} --- Slow, conscious, effortful (System 2)
\item $L_3$: \textbf{Meta-Cognitive} --- Thinking about thinking
\end{itemize}

\textbf{Key Insight:}
Different biases and interventions apply at different levels.

\textbf{Cross-References:}
AAA (CORE-WHO), V (CORE-WHEN), AU (CORE-AWARE), IE (CORE-EIT)

\end{tcolorbox}

\section{The Fundamental Question}
\label{sec:hier-fundamental}

\begin{tcolorbox}[colback=red!5!white, colframe=red!75!black,
    title=The Fundamental Question]
\textbf{At what level of the decision hierarchy is this choice being made?
Which cognitive system is engaged?}
\end{tcolorbox}

This is distinct from CORE-WHO (welfare hierarchy):
\begin{itemize}
\item WHO: \textit{Whose} utility matters (individual, group, society)
\item HIERARCHY: \textit{How} the decision is processed (automatic vs deliberative)
\end{itemize}

\section{Core Theory: The Four Decision Levels}
\label{sec:hier-levels}

\begin{table}[htbp]
\centering
\caption{The Four Decision Levels}
\label{tab:hier-levels}
\renewcommand{\arraystretch}{1.4}
\begin{tabular}{@{}clccp{3.5cm}@{}}
\toprule
\textbf{Level} & \textbf{Name} & \textbf{Speed} & \textbf{Effort} & \textbf{Characteristic} \\
\midrule
$L_0$ & Automatic & Fast & Low & Unconscious, habitual \\
$L_1$ & Heuristic & Medium & Low-Med & Rule-of-thumb, satisficing \\
$L_2$ & Deliberative & Slow & High & Conscious, analytical \\
$L_3$ & Meta-Cognitive & Variable & High & Reflection on thinking \\
\bottomrule
\end{tabular}
\end{table}

\subsection{Level Characteristics}

\subsubsection{$L_0$: Automatic Processing}
\begin{itemize}[nosep]
\item Corresponds to Kahneman's ``System 1''
\item Fast, parallel, effortless
\item Subject to: anchoring, framing, availability
\item Intervention approach: defaults, choice architecture
\end{itemize}

\subsubsection{$L_1$: Heuristic Processing}
\begin{itemize}[nosep]
\item Uses simple rules and shortcuts
\item ``Good enough'' rather than optimal
\item Subject to: representativeness, recognition heuristic
\item Intervention approach: simplification, salient cues
\end{itemize}

\subsubsection{$L_2$: Deliberative Processing}
\begin{itemize}[nosep]
\item Corresponds to Kahneman's ``System 2''
\item Slow, serial, effortful
\item Subject to: cognitive depletion, overconfidence
\item Intervention approach: information, decision aids
\end{itemize}

\subsubsection{$L_3$: Meta-Cognitive Processing}
\begin{itemize}[nosep]
\item Thinking about thinking
\item Monitoring own decision processes
\item Subject to: illusion of introspection
\item Intervention approach: debiasing training, reflection prompts
\end{itemize}

\section{Level Determinants}
\label{sec:hier-determinants}

The active decision level depends on:

\begin{equation}
L = f(\text{stakes}, \text{time}, \text{load}, \text{novelty}, \text{expertise})
\end{equation}

\begin{table}[htbp]
\centering
\caption{Factors Affecting Decision Level}
\label{tab:hier-determinants}
\renewcommand{\arraystretch}{1.3}
\begin{tabular}{@{}lcc@{}}
\toprule
\textbf{Factor} & \textbf{Low} & \textbf{High} \\
\midrule
Stakes & $L_0, L_1$ & $L_2, L_3$ \\
Time pressure & $L_0, L_1$ & $L_1, L_2$ \\
Cognitive load & $L_0$ & $L_0, L_1$ \\
Novelty & $L_0$ (if habit) & $L_2$ \\
Expertise & $L_0$ (intuition) & $L_2$ (if novel) \\
\bottomrule
\end{tabular}
\end{table}

\section{Axiom System}
\label{sec:hier-axioms}

\begin{tcolorbox}[colback=gray!5!white, colframe=gray!75!black,
    title=Axiom HIE-1: Level-Bias Correspondence]
\textbf{Statement:}
\begin{equation}
\text{Bias}_b \text{ active} \Leftrightarrow L \in \mathcal{L}_b
\end{equation}

\textbf{Interpretation:} Each bias has a set of levels where it operates.
Anchoring operates at $L_0$; overconfidence at $L_2$.
\end{tcolorbox}

\begin{tcolorbox}[colback=gray!5!white, colframe=gray!75!black,
    title=Axiom HIE-2: Level Shifting]
\textbf{Statement:}
\begin{equation}
L_i \to L_j \text{ possible via intervention or context change}
\end{equation}

\textbf{Interpretation:} Interventions can shift processing level.
Prompts can move from $L_0$ to $L_2$; time pressure moves from $L_2$ to $L_0$.
\end{tcolorbox}

\begin{tcolorbox}[colback=gray!5!white, colframe=gray!75!black,
    title=Axiom HIE-3: Default Level]
\textbf{Statement:}
\begin{equation}
L_{\text{default}} = L_0 \quad \text{(automatic processing is default)}
\end{equation}

\textbf{Interpretation:} Without intervention, decisions default to automatic processing.
Deliberation requires activation energy.
\end{tcolorbox}

\begin{tcolorbox}[colback=gray!5!white, colframe=gray!75!black,
    title=Axiom HIE-4: Level-Intervention Match]
\textbf{Statement:}
\begin{equation}
\text{Effectiveness}(I) \propto \text{Match}(I.\text{target\_level}, L_{\text{actual}})
\end{equation}

\textbf{Interpretation:} Interventions designed for one level underperform when
applied at another level.
\end{tcolorbox}

\section{Number of $L_2$ Decisions ($N_{L2}$)}
\label{sec:hier-NL2}

\begin{definition}[$N_{L2}$ Capacity]
$N_{L2}$ represents the number of deliberative decisions an individual can make
per day before cognitive depletion reduces all subsequent decisions to $L_0$ or $L_1$.
\end{definition}

\begin{equation}
N_{L2} \approx 5-10 \text{ major decisions/day}
\end{equation}

\textbf{Implication:} Important decisions should be scheduled early;
routine decisions should be automated or defaulted.

\section{Integration with Other COREs}

\begin{tcolorbox}[colback=yellow!5!white, colframe=yellow!75!black,
    title=Integration: CORE-HIERARCHY $\leftrightarrow$ Other COREs]
\textbf{With CORE-WHO (AAA):}
Decision hierarchy is orthogonal to welfare hierarchy. A group decision can be
automatic ($L_0$) or deliberative ($L_2$).

\textbf{With CORE-WHEN (V):}
Context determines active level: time pressure → $L_0$; high stakes → $L_2$.

\textbf{With CORE-AWARE (AU):}
Higher levels ($L_2$, $L_3$) typically involve higher awareness.

\textbf{With CORE-EIT (IE):}
Intervention type must match decision level for effectiveness.
\end{tcolorbox}

\section{Summary}

\begin{tcolorbox}[colback=gray!5!white, colframe=gray!75!black,
    title=Appendix HI Summary: CORE-HIERARCHY]

\textbf{Core Contribution:}
Formalizes the four-level decision hierarchy that determines which biases
apply and which interventions are effective.

\textbf{Key Insight:}
\begin{itemize}[nosep]
\item $L_0$ (Automatic): Defaults, choice architecture
\item $L_1$ (Heuristic): Simplification, salient cues
\item $L_2$ (Deliberative): Information, decision aids
\item $L_3$ (Meta-Cognitive): Debiasing, reflection prompts
\end{itemize}

\textbf{Distinction from CORE-WHO:}
WHO = whose utility; HIERARCHY = how decision is processed.

\end{tcolorbox}

\vspace{1em}
\hrule
\vspace{0.5em}

\noindent\textit{Cross-references:}
Appendix AAA (CORE-WHO), V (CORE-WHEN), AU (CORE-AWARE), IE (CORE-EIT).

\noindent\textit{Master Bibliography:} \texttt{bcm\_master.bib}

\nocite{bcm_master}

% =============================================================================
% END OF APPENDIX HI
% =============================================================================
